\begin{textsample}{POS Dim 2 – vem\_ed\_21 – Score 1.00 – vem\_ed\_21/t103.txt}  \label{ex:f2_pos_171}
Leitores Aos Osvaldo Succi Jr.
Coordenador PCIs Este número 21 de VEm está dedicado à cobertura da principal conferência internacional sobre Intercâmbios Virtuais.
A IVEC 2023 ocorreu de 30 de outubro a 1 de novembro de 2023, pela primeira vez no Hemisfério Sul, no Brasil (em São Paulo).
Contou com 496 participantes, presenciais e virtuais, de 41 países, e foi organizada pela Associação Brasileira de Educação Internacional (Faubai) e pela Unesp.
Os PCIs / Cesu demonstram o vanguardismo de suas ações no Brasil e no mundo e, na conferência IVEC, isso \textbf{ficou} claro por alguns motivos: a quantidade de trabalhos aceitos (duas oficinas e três comunicações orais); a conquista de bolsas para cobertura de gastos com inscrição e estadia, oferecidas pela Stevens Initiative, patrocinadora “Platinum” do evento, por três membros da equipe (Ana Carolina Freschi, Regiane Souza Camargo Moreira e eu); a participação como moderadores em quatro sessões de comunicação oral; o convite da Unesp e da Faubai à professora Patrícia Sales Patrício, jornalista responsável pela comunicação dos PCIs, para realizar a cobertura do evento; a participação deste coordenador da equipe de PCIs no painel de encerramento da conferência, “What is the future of COIL / virtual exchange?” (Qual o futuro do COIL / Intercâmbio Virtual).
Nas páginas a seguir, você encontra uma síntese dessas produções.
Boa leitura!

% matched lemmas: ficar
\end{textsample}
